\clearpage
\phantomsection% 
\addcontentsline{toc}{chapter}{ABSTRACT}
\thispagestyle{fancy}

\begin{center}
    \large \bfseries \MakeUppercase{Abstract}\\
    \normalsize \normalfont {\thetitleEN}\\
    \normalsize \normalfont {\theauthor}\\
    \bigskip
    
    \normalsize \normalfont \justifying \singlespacing
    The prevalence of stunting in Indonesia is a communal health problem influenced by various complex and non-linear determinant factors. This study aims to identify the main determinants of stunting at the provincial level and extract their critical threshold points using the aggregated data of the 2024 Indonesian Nutritional Status Survey (SSGI). The modeling methodology was built using the Random Forest Regression algorithm, while the model's interpretability to unlock its black-box nature was analyzed through an Explainable AI approach using the SHapley Additive exPlanations (SHAP) method. Evaluation results indicate that the model has excellent predictive performance, evidenced by a coefficient of determination ($R^2$) of 0.9090, Mean Absolute Error (MAE) of 1.40, Root Mean Squared Error (RMSE) of 1.77, and Mean Absolute Percentage Error (MAPE) of 7.22\%. Based on the SHAP feature importance visualization, the most dominating determinants are the lack of daily nutritional standards for toddlers in the form of the absence of Minimum Dietary Diversity (MDD) and Minimum Milk Feeding Frequency (MMFF), a high percentage of at-risk maternal age, infrastructure access (risky sanitation), and the lack of Supplementary Feeding (PMT). Furthermore, the Dependence Plot analysis successfully revealed communal mathematical tolerance thresholds, such as the absence of MDD at 55\%, the absence of MMFF at 22\%, at-risk maternal age at 14\%, and risky sanitation facilities at 10\%. A drastic surge in the risk of stunting prevalence occurs when a region exceeds these percentage thresholds. In conclusion, the predictive model and the resulting visual explanations have proven effective in precisely mapping the root causes of malnutrition, making it a viable baseline for evaluating national stunting intervention policies.\par
    
    \vspace{20pt}
    \raggedright\textbf{Keywords: Stunting, Random Forest, SHAP, SSGI, Determinant}
    
    \vfill
    
\end{center}
\clearpage